% TOP_PROBLEM: {"id": 2985, "title": "Kirby Problem 4.109", "category_id": 7, "category": {"id": 7, "name": "topology", "display_name": "Topology", "description": "Properties preserved under continuous deformations.", "slug": "topology", "order_index": 7, "created_at": "2026-07-31T15:26:25.671Z"}, "status": "open"}
% FIRST_POSTED: null
% ENTRY_KIND: statement_only
% Imported from: batch13.tex; UnsolvedMath ID 2985
\documentclass[11pt]{article}
\usepackage[margin=25mm]{geometry}
\usepackage{fontspec}
\setmainfont{Latin Modern Roman}
\usepackage{amsmath,amssymb,mathtools}
\usepackage{unicode-math}
\setmathfont{Latin Modern Math}
\usepackage{xurl,hyperref}
\hypersetup{colorlinks=true,urlcolor=blue}
\setcounter{secnumdepth}{0}
\setlength{\parindent}{0pt}
\setlength{\parskip}{5pt}
\setlength{\emergencystretch}{3em}
\newcommand{\sref}[2]{\hyperlink{source:#1:#2}{[S#2]}}
\newcommand{\cataloguescope}[1]{\paragraph{Recorded target.}#1}
\begin{document}
% UnsolvedMath ID 2985; source catalogue code KP-4.109
\setcounter{section}{2984}
\section{Kirby Problem 4.109}\label{2985}
{\small\sffamily UnsolvedMath ID 2985 · Topology}
\subsection{Definitions and mathematical statement}

\paragraph{Shared notation.}$\mathbb N=\{1,2,\ldots\}$, and $\mathbb Z,\mathbb Q,\mathbb R,\mathbb C$ denote the integers, rationals, reals and complex numbers. Any other conventions are specified locally.


Let $(X,\omega)$ be a closed symplectic four-manifold, meaning that $\omega$ is a closed nondegenerate real two-form. Let $\Sigma\subset X$ be a closed embedded symplectic surface: the restriction of $\omega$ to each oriented tangent plane of $\Sigma$ is positive. Suppose that, in real homology,
\[
[\Sigma]=\operatorname{PD}(k[\omega])\in H_2(X;\mathbb R)
\]
for an integer $k$, where $\operatorname{PD}$ is Poincaré duality and the integral fundamental class is viewed in real homology. A Weinstein domain is a compact manifold $F$ with boundary, a one-form $\lambda$ such that $d\lambda$ is symplectic and its Liouville vector field $V$ points outward along the boundary, and a Morse function for which $V$ is gradient-like. Here $\iota_Vd\lambda=\lambda$; gradient-like means that the function increases strictly along $V$ away from its critical points, with the usual local gradient behavior there. In the complement convention of Giroux, compatibility with the given form means that there is a diffeomorphism
\[
 q:\operatorname{int}(F)\longrightarrow X\setminus\Sigma,
 \qquad q^*(\omega|_{X\setminus\Sigma})=d\lambda|_{\operatorname{int}(F)}.
\]
The symplectic complement is the interior of a Weinstein domain, rather than a completion obtained by attaching an infinite-volume end. \sref{2985}{3}
\paragraph{Statement recorded in the supplied source.}
For every such symplectic surface $\Sigma$, does its complement $X\setminus\Sigma$ admit a Weinstein structure compatible with the restriction of the given form $\omega$? The question concerns each given representative of the indicated Poincaré-dual class, rather than the existence of some suitably chosen representative. \sref{2985}{2}
\subsection{Short English statement}
Must the complement of every symplectic surface in an integral multiple of the symplectic class carry a Weinstein structure whose symplectic form is the original restricted form?
\subsection{Sources}
\begin{description}
\item[\hypertarget{source:2985:1}{S1}] ULAM.AI and the UnsolvedMath Contributors, UnsolvedMath: A Curated Collection of Open Mathematics Problems, record 2985 (KP-4.109). CC BY 4.0. \url{https://huggingface.co/datasets/ulamai/UnsolvedMath}.
\item[\hypertarget{source:2985:2}{S2}] R. İnanç Baykur, Robion C. Kirby and Daniel Ruberman, K3—A New Problem List in Low-Dimensional Topology, Mathematical Surveys and Monographs 295 (2026), author version, Problem 4.109 and Remarks, PDF page 281. \url{https://math.berkeley.edu/sites/default/files/surv-295-ruberman-watermarked-author-pdf.pdf}.
\item[\hypertarget{source:2985:3}{S3}] Emmanuel Giroux, Remarks on Donaldson's symplectic submanifolds, arXiv:1803.05929v1 (2018), pages 2--3, Weinstein-domain definition and Theorem 2. \url{https://arxiv.org/pdf/1803.05929v1}.
\end{description}
\label{end:2985}
\end{document}
